\chapter{Conclusions and Further Work}
\label{ch:conclusion}

% GUIDANCE (chapter target about 800 words). Summarise with acumen: what was asked, what was built,
% what was found, what it means, in that order — then further work that follows FROM the evaluation,
% not a wish list.

\section{Conclusions}
\label{sec:conc}
% GUIDANCE: Map back to the objectives in Section~\ref{sec:intro-aims}, one sentence each, with the
% deciding number as a macro (e.g. \FullEmoSignObs\ vs \FullEmoSignNull; \DiffSignPp; \McDiffFamPp;
% \ConsFam\ vs \ConsFamNull). End with the one-sentence answer to the question in
% Section~\ref{sec:intro-question}, bounded exactly as Section~\ref{sec:disc-not-claimed} bounds it.

\section{Further work}
\label{sec:further}
% GUIDANCE: In priority order, each tied to a limitation or an audit result:
%   1. A human-rated study of the drawings (the study this pipeline was built to precede), rating the
%      same ladder so machine and human agreement can be compared directly; \cite{lomas2024improved}
%      is the design to adapt.
%   2. A people-aware description template as the default (no expression clause when the describer
%      reports no people), with the demonstrator kept in parity; the ablation is its evidence.
%   3. A second describer (a sketch- or drawing-specific model was designed but not executed) and a
%      test--retest estimate for both readers, so the consistency result can be decomposed.
%   4. NRC-VAD per-family means in place of hand-set centroids \cite{mohammad2018obtaining}.
%   5. A richer family scheme (or the 27 categories of \cite{cowen2017self}) so low-arousal positive
%      states do not fall to neutral.
%   6. Extending the single-coder check to a second coder with an agreement statistic.
